\begin{frame}[label=q5a2]{5a.2 · DataFrame: año, sexo y filtro agregado}
  \small
  \pregunta{¿En qué años el total reportado de mujeres supera las 5\,000 personas?}
  \entrada{93\,831 filas. \texttt{anio}, \texttt{sexo}, \texttt{cantidad}; salida filtrada a F.}
  \flujo{Todos los registros}{\texttt{groupBy(anio,sexo)}}{\texttt{SUM(cantidad)}\\Sexo F y total > 5000}{6 años\\Orden ascendente}
  \begin{columns}[T,onlytextwidth]
    \column{0.55\textwidth}
    \textbf{Código Scala original (fragmento)}
    \codigo{sql/05a_2_anio_sexo_mujeres}{10}{16}
    \column{0.42\textwidth}
    \textbf{Salida}\\[5pt]
    {\footnotesize
\begin{tabular}{@{}lll@{}}
\toprule
Año & Personas F & Registros\\
\midrule
2019 & 16\,451 & 11\,037\\
2020 & 11\,828 & 8\,723\\
2024 & 11\,275 & 8\,398\\
\bottomrule
\end{tabular}
    }
  \end{columns}
  \respuesta{Los seis años cumplen el umbral. En 2024 se contabilizan 11\,275 mujeres en 8\,398 filas, dos medidas distintas.}
  \origen{sql}{05a_2_anio_sexo_mujeres}
\end{frame}
